% Preamble for the study guide. Built with tectonic (XeTeX engine).
\usepackage{amsmath}
\usepackage{mathtools}
\usepackage[mono=false]{libertinus}
\usepackage{fontspec}
\setmonofont{DejaVuSansMono.ttf}[
  Scale=0.80,
  HyphenChar=None,
  BoldFont=DejaVuSansMono-Bold.ttf,
  ItalicFont=DejaVuSansMono-Oblique.ttf,
  BoldItalicFont=DejaVuSansMono-BoldOblique.ttf]
\usepackage{polyglossia}
\setdefaultlanguage{swedish}
\setotherlanguage{english}
\usepackage[a4paper,left=2.5cm,right=2.5cm,top=2.3cm,bottom=2.6cm,footskip=1.1cm]{geometry}
\usepackage{microtype}
\usepackage{xcolor}
\usepackage{graphicx}
\usepackage{booktabs}
\usepackage{tabularx}
\usepackage{array}
\usepackage{multicol}
\usepackage{enumitem}
\usepackage{siunitx}
\usepackage{float}
\usepackage{listings}
\usepackage{tikz}
\usetikzlibrary{arrows.meta,positioning,fit,calc}
\usepackage[most]{tcolorbox}
\usepackage[hidelinks,unicode,bookmarksnumbered,bookmarksopen,bookmarksopenlevel=1]{hyperref}
\usepackage{bookmark}

\graphicspath{{../figures/}}

% ---------------------------------------------------------------- colours
\definecolor{accent}{HTML}{1F4E79}
\definecolor{boxblue}{HTML}{1F5FA8}
\definecolor{boxred}{HTML}{B5402F}
\definecolor{boxgreen}{HTML}{217A52}
\definecolor{boxgray}{HTML}{5F6670}
\definecolor{codebg}{HTML}{F5F6F8}
\definecolor{codeframe}{HTML}{D4D8DE}
\definecolor{codegray}{HTML}{6A737D}
\definecolor{codeblue}{HTML}{1F4E79}
\definecolor{codegreen}{HTML}{2D6A4F}

% ---------------------------------------------------------------- KOMA
\KOMAoptions{parskip=half-,headings=small,captions=tableheading,toc=flat}
\DeclareTOCStyleEntries[raggedentrytext=true]{tocline}{section,subsection}
\setkomafont{disposition}{\sffamily\bfseries\color{accent}}
\setkomafont{title}{\sffamily\bfseries\color{accent}}
\setkomafont{subtitle}{\sffamily\color{black!70}}
\setkomafont{caption}{\small}
\setkomafont{captionlabel}{\small\sffamily\bfseries\color{accent}}
\setcapindent{0pt}
\addtokomafont{descriptionlabel}{\rmfamily\bfseries}
\RedeclareSectionCommand[beforeskip=1.6\baselineskip,afterskip=0.6\baselineskip]{section}
\RedeclareSectionCommand[beforeskip=1.0\baselineskip,afterskip=0.3\baselineskip]{subsection}
\RedeclareSectionCommand[beforeskip=0.7\baselineskip,afterskip=0.15\baselineskip]{subsubsection}
\addtokomafont{section}{\Large}
% Little stretch between paragraphs: with the default (plus half a line) a ragged-bottom page
% can be cut short at any heading, which leaves large gaps at the foot of the page.
\setparsizes{0pt}{0.5\baselineskip plus 0.6pt}{0pt plus 1fil}
\setlength{\textfloatsep}{14pt plus 2pt minus 2pt}
\setlength{\intextsep}{12pt plus 2pt minus 2pt}
\setlength{\floatsep}{12pt plus 2pt minus 2pt}
\setcounter{secnumdepth}{2}
\setcounter{tocdepth}{2}
\numberwithin{equation}{section}
\numberwithin{figure}{section}
\numberwithin{table}{section}
\renewcommand{\floatpagefraction}{0.8}
\renewcommand{\topfraction}{0.9}
\renewcommand{\bottomfraction}{0.7}
\renewcommand{\textfraction}{0.08}
\setlength{\emergencystretch}{2em}
\setlist{itemsep=0.15em,topsep=0.3em,parsep=0pt,leftmargin=1.6em}
\setlist[description]{leftmargin=1.6em,labelindent=0pt}

% ---------------------------------------------------------------- decimal comma
% Same behaviour as the icomma package (a comma followed by a space is punctuation,
% otherwise a decimal comma), but set up after unicode-math so that the comma glyph
% comes from the math font in use.
\makeatletter
\AtBeginDocument{%
  \Umathcharnumdef\mathcomma=\Umathcodenum`\,\relax
  \mathcode`\,="8000\relax}
\begingroup\catcode`\,=\active
  \gdef,{\futurelet\@let@token\sm@rtcomma}
\endgroup
\def\sm@rtcomma{%
  \ifx\@let@token\@sptoken\else\ifx\@let@token\space\else\mathord\fi\fi\mathcomma}
\makeatother

% ---------------------------------------------------------------- numbers
\sisetup{output-decimal-marker={,},group-separator={\,},group-minimum-digits=4,
         detect-all,mode=match}
\newcommand{\pct}{\,\%}
\newcommand{\kr}{\,kr}
\newcommand{\x}{\ensuremath{\times}}          % "19,4\x" = 19,4 times

% ---------------------------------------------------------------- text macros
\newcommand{\eng}[1]{(\textenglish{\emph{#1}})}
\newcommand{\en}[1]{\textenglish{\emph{#1}}}
\newcommand{\kod}[1]{\texttt{\detokenize{#1}}}
\newcommand{\flag}[1]{\texttt{-{}-\detokenize{#1}}}   % command line flag
\newcommand{\EV}{\mathrm{EV}}
\newcommand{\TV}{\mathrm{TV}}
\newcommand{\FCFF}{\mathrm{FCFF}}
\newcommand{\FCFE}{\mathrm{FCFE}}
\newcommand{\EBITDA}{\mathrm{EBITDA}}
\newcommand{\EBIT}{\mathrm{EBIT}}
\newcommand{\NOPAT}{\mathrm{NOPAT}}
\newcommand{\WACC}{\mathrm{WACC}}
\newcommand{\ERP}{\mathrm{ERP}}
\newcommand{\NWC}{\mathrm{NWC}}
\newcommand{\PV}{\mathrm{PV}}
\newcommand{\IRR}{\mathrm{IRR}}
\newcommand{\MOIC}{\mathrm{MOIC}}
\newcommand{\Cov}{\operatorname{Cov}}
\newcommand{\Var}{\operatorname{Var}}
\newcommand{\E}{\operatorname{E}}
\newcommand{\T}{^{\mathsf{T}}}

% ---------------------------------------------------------------- callouts
\newtcolorbox{callout}[2]{enhanced,breakable=false,colback=#1!5!white,colframe=#1,
  boxrule=0pt,leftrule=3pt,arc=0pt,outer arc=0pt,left=8pt,right=8pt,top=5pt,bottom=5pt,
  before skip=0.8\baselineskip,after skip=0.8\baselineskip,
  fonttitle=\sffamily\bfseries\small,coltitle=#1,colbacktitle=#1!5!white,
  title={#2},attach title to upper={\quad},parbox=false}
\newenvironment{komihag}{\begin{callout}{boxblue}{Kom ihåg}}{\end{callout}}
\newenvironment{vanligtfel}{\begin{callout}{boxred}{Vanligt fel}}{\end{callout}}
\newenvironment{intervjutips}{\begin{callout}{boxgreen}{Intervjutips}}{\end{callout}}
\newenvironment{imotorn}{\begin{callout}{boxgray}{Så gör motorn}}{\end{callout}}

% ---------------------------------------------------------------- listings
\lstdefinestyle{base}{basicstyle=\ttfamily\small,columns=fullflexible,keepspaces=true,
  backgroundcolor=\color{codebg},frame=single,rulecolor=\color{codeframe},framerule=0.4pt,
  xleftmargin=6pt,xrightmargin=6pt,framesep=5pt,aboveskip=0.7\baselineskip,
  belowskip=0.5\baselineskip,showstringspaces=false,breaklines=true,breakatwhitespace=true,
  upquote=true,literate={-}{{-}}1}
\lstdefinestyle{py}{style=base,language=Python,keywordstyle=\color{codeblue}\bfseries,
  commentstyle=\color{codegray}\itshape,stringstyle=\color{codegreen}}
\lstdefinestyle{sh}{style=base,language=bash,keywordstyle=\color{black},
  commentstyle=\color{codegray}\itshape,stringstyle=\color{black}}
\lstdefinestyle{txt}{style=base,language={}}
\lstset{style=py}
\AtBeginDocument{\renewcommand{\lstlistingname}{Kodexempel}}

% ---------------------------------------------------------------- Q&A and exercises
\newcounter{fraga}
\newenvironment{fraga}[1]{\par\addvspace{0.9\baselineskip}\refstepcounter{fraga}%
  \noindent\begin{minipage}{\linewidth}\setlength{\parskip}{0.4\baselineskip plus 0.6pt}%
  \noindent\textsf{\textbf{\color{accent}Q\thefraga.}}\enspace\textenglish{\textbf{#1}}\par
  \noindent\ignorespaces}{\end{minipage}\par}
\newcommand{\fraser}[1]{\par\noindent{\small\textsf{\textbf{Säg på engelska:}}
  \textenglish{\itshape #1}}\par}
\newcounter{ovning}
\newenvironment{ovning}[1][]{\par\refstepcounter{ovning}\addvspace{0.8\baselineskip}%
  \noindent\textsf{\textbf{\color{accent}Övning \theovning}}%
  \if\relax\detokenize{#1}\relax\else\ \textsf{(#1)}\fi\textsf{\textbf{.}}\enspace}{\par}
\newenvironment{losning}{\par\vspace{0.2\baselineskip}\noindent\textsf{\textbf{Lösning.}}\enspace}{\par}
